\section{The Three Softmaxes}
\label{sec:three-softmaxes}

There are three places in the transformer where softmax or sigmoid appears. They do
three completely different jobs. Conflating them is the single most common source of
confusion about what the transformer is computing.

\begin{center}
\begin{tabular}{llll}
\toprule
& Job & Input & Semantics \\
\midrule
Softmax 1 & Routing & $Q \cdot K$ scores & Differentiable argmax \\
Softmax 2 & Inference & Neighbor beliefs & Bayesian posterior \\
Softmax 3 & Generation & Final hidden state & Token distribution \\
\bottomrule
\end{tabular}
\end{center}

\subsection*{Softmax 1: Attention (Routing)}

The first softmax appears inside every attention head, applied to the query-key dot
products. Its job is routing: which token should I look at? This is a differentiable
argmax --- at high temperature it concentrates all weight on the token with the
highest $Q \cdot K$ score.

In the BP construction, the attention softmax fetches a neighbor's belief by index
matching. The query encodes ``I am looking for the token at index $i$''; the key
encodes ``I am the token at index $i$''; the dot product peaks at the correct match;
softmax concentrates the weight there. The result is a lookup table.

This softmax is \emph{not} doing inference. It is not computing a probability over
hypotheses. It is selecting which piece of evidence to retrieve. Calling it a
``probability distribution over tokens'' is technically correct but misleading ---
the weights are routing coefficients, not posterior beliefs.

\subsection*{Softmax 2: The FFN Sigmoid (Inference)}

The second softmax is the sigmoid activation in the FFN. Its job is Bayesian
inference: what is my updated belief given my neighbors' beliefs?
\[
\updateBelief(m_0, m_1) = \sigma\!\left(\logit(m_0) + \logit(m_1)\right).
\]
This is where the actual Bayesian computation happens. The sigmoid is the exact
inverse of logit, converting the log-odds sum back to a probability. This is the
weight-of-evidence combination of Turing and Good~\cite{good1950probability},
applied to two binary evidence sources.

This softmax \emph{is} computing a posterior. The output is a belief in $(0,1)$
representing the updated probability that the current proposition is true, given
the gathered evidence. This is the inference step. Everything else in the
transformer is infrastructure for getting the right inputs to this computation.

\subsection*{Softmax 3: Output (Generation)}

The third softmax appears at the output layer, applied to the logits over the
vocabulary. Its job is generation: which token comes next? In a standard language
model, the logits are computed by a matrix multiply over the final hidden state ---
pattern-matched association rather than computed inference.

In a grounded QBBN transformer, the output would be different: the logits would come
from computed BP posteriors, representing the marginal probability of each
proposition being true given the evidence. The softmax operation is the same; what
changes is what it is applied to.

The QBBN does not replace the softmax. It replaces the logits. The architecture was
always capable of computing exact posteriors; it just needed the right weights to
feed into the output softmax. Standard language model training produces weights that
implement pattern-matched association. BP training would produce weights that
implement principled inference.

\subsection*{Why the Distinction Matters}

Conflating Softmax~1 and Softmax~2 leads to the mistaken conclusion that the
attention weights are the inference. They are not. The attention weights determine
\emph{which} evidence to gather. The FFN sigmoid determines \emph{what to conclude}
from that evidence. Routing and inference are separate computations, performed by
separate components, in a strict temporal order enforced by the architecture.

Conflating Softmax~2 and Softmax~3 leads to the mistaken conclusion that the output
distribution is a posterior over propositions. In an ungrounded model, it is not ---
it is a distribution over next tokens learned by association. The two converge only
in a grounded model where every output token corresponds to a declared proposition
with a known prior.